\honest{Schools Proliferating Without Evidence}{Eliezer Yudkowsky}{2009}

Robyn Dawes, one of the few who try to bring the research on judgment into real life, wrote \textsc{House of Cards}. I quote its first chapter: virtually all the research has found that psychotherapists' ``claims to superior intuitive insight, understanding, and skill as therapists are simply invalid''. \nb{The book could not be checked here.}

Remember the Rorschach? Hundreds of experiments have looked for evidence that it works, and the answer, I say, is simply ``No.'' Yet it is still in use. \nb{The most critical review before the post, by Lilienfeld, Wood and Garb in 2000, found support for ``a small number of indexes derived from the Rorschach and TAT'', a related test, and none for most of them.} That tells you what sort of field this is.

The rest of the evidence agrees. Patients who see therapists get better faster than those who do nothing. But the schools of therapy do not differ, and experience brings no gain. \nb{The claims about schools and about experience had the support of major reviews, and the second has held up in later work. The finding on schools, called the dodo bird verdict, was qualified by evidence for cognitive-behavioral therapy, which the post's added line recalls.} Talking to ``a randomly selected college professor from another field'' works as well as seeing a therapist. ``It's just talking to anyone that helps you get better, apparently.'' \nb{The claim matches one study, by Strupp and Hadley in 1979. Its professors were ``chosen for their ability to form understanding relationships'', not selected at random, and its patients were male college students.}

``In the entire absence of the slightest experimental evidence for their effectiveness,'' therapists were licensed, heard in court, accredited and paid by insurers. \nb{Two paragraphs earlier the post grants that patients in therapy get better faster than those who do nothing.} Schools multiplied. ``So far as I can discern,'' prestige came from founding a school on charisma and good stories, as Freud did, and this ``probably'' explains why psychotherapy exists at all. I give no evidence.

Happiness research, by contrast, began when people found measures of happiness that agree with each other. \nb{They do: in 1993 Sandvik, Diener and Seidlitz found that the usual self-report measures ``correlated highly with alternative measures''.} So any organized practice needs ``a control group, an experimental group, and statistics''. \nb{By the post's own account, psychotherapy research had these; its findings are ``the experimental results on the field''. What failed was acting on them, and the prescription does not cover that.}

In an addition I say that ``Dawes wrote in the 80s'', that matters may have improved, and that I remember hearing of evidence for cognitive-behavioral therapy. \nb{\textsc{House of Cards} appeared in 1994.}
